\documentclass[12pt, a4paper]{report}

% --- Packages ---
\usepackage[utf8]{inputenc}
\usepackage[T1]{fontenc}
\usepackage{geometry}
\usepackage{graphicx}
\usepackage{hyperref}
\usepackage{xcolor}
\usepackage{titlesec}
\usepackage{fancyhdr}
\usepackage{booktabs}
\usepackage{enumitem}
\usepackage{setspace}
\usepackage{tikz}
\usepackage{float}
\usepackage{subcaption}
\usepackage{times} % Research paper friendly font

\usetikzlibrary{shapes.geometric, arrows, positioning, shadows}

% --- Configuration ---
\geometry{left=1.5cm, right=1.5cm, top=1.5cm, bottom=2.5cm}
\hypersetup{
    colorlinks=true,
    linkcolor=black,
    filecolor=magenta,      
    urlcolor=blue,
    pdftitle={MediTrack},
    pdfpagemode=FullScreen,
}
\setlength{\parskip}{0.5em}
\setlength{\parindent}{0pt}
\onehalfspacing



% --- Custom Colors ---
\definecolor{primaryBlue}{RGB}{0, 51, 102}
\titleformat{\chapter}[display]
  {\normalfont\huge\bfseries\color{primaryBlue}}{\chaptertitlename\ \thechapter}{20pt}{\Huge}
\titleformat{\section}
  {\normalfont\Large\bfseries\color{primaryBlue}}{\thesection}{1em}{}

% --- TikZ Style for DFD ---
\tikzstyle{entity} = [rectangle, draw=black, fill=gray!20, text width=2.5cm, text centered, rounded corners, minimum height=1.5cm]
\tikzstyle{process} = [circle, draw=black, fill=white, text width=3cm, text centered, minimum height=3cm, drop shadow]
\tikzstyle{arrow} = [thick,->,>=stealth]

% =========================================
% DOCUMENT BEGINS
% =========================================
\begin{document}

% --- Title Page ---
\begin{titlepage}
    \begin{center}
        % 1. University Logo at the VERY TOP
        \includegraphics[width=0.25\textwidth]{Images/CoU Logo.png} % Replace with actual logo file
        \vspace{.60cm} % Space after logo
        
        % 2. Project Title
        \Huge
        \textbf{\textcolor{primaryBlue}{MediTrack}}
        
        \vspace{0.2cm}
        \LARGE
        A Progressive Web Application (PWA) for \\ Family Health Management
        
        \vspace{1.5cm} % Space before course info
        
        % 3. Course Information
        \textbf{Course Title:} Web Engineering Lab \\
        \textbf{Course Code:} CSE 4103 % Replace with actual code
        \vspace{0.2cm}
        
        \vfill % This pushes the following content to the bottom half
        
        % 4. Submission Information
        \Large
        \textbf{Submitted By:} \\
        Md. Al-Amin \\
        ID: 12108036 \\
        Department of Computer Science and Engineering
        
        \vspace{1cm}
        
        \textbf{Submitted To:} \\
        Partha Chakraborty \\
        Assistant Professor \\
        Department of Computer Science and Engineering
        
        \vspace{1.5cm}
        
        % 5. Date
        \large
        Date: December 4, 2025
        
    \end{center}
\end{titlepage}

% --- Table of Contents ---
\tableofcontents
\newpage

% =========================================
% CHAPTER 1: INTRODUCTION
% =========================================
\chapter{Introduction}

Family health management faces significant challenges in Bangladesh, including inadequate infrastructure, limited access to quality care in rural areas, and rising prevalence of chronic diseases. Poor medication adherence and fragmented health information across multiple providers contribute to suboptimal health outcomes. Traditional paper-based systems and disconnected digital applications fail to provide unified family health management solutions.

\section{Project Overview: MediTrack}
MediTrack is a full-stack Progressive Web Application (PWA) built with React, Vite, Express.js, and Prisma, providing:

\begin{itemize}
    \item \textbf{Family Members Management:} Users can create and manage profiles for all family members, tracking their health information, medical history, and relationships.
    \item \textbf{Medication Database:} Comprehensive medication tracking with scheduling, dosage management, and adherence monitoring.
    \item \textbf{Appointment Management:} Centralized scheduling system for tracking medical appointments across family members.
    \item \textbf{Health Reports:} Secure storage and organization of medical documents, test results, and health reports.
    \item \textbf{Analytics and Insights:} Data visualization and health trend analysis using Recharts for informed decision-making.
    \item \textbf{Purchase List:} Medication purchase tracking and cost management features.
\end{itemize}

MediTrack demonstrates practical implementation of modern React architecture, RESTful API design, PWA principles, secure authentication, and responsive design.

% =========================================
% CHAPTER 2: LITERATURE REVIEW
% =========================================
\chapter{Literature Review}

\section{Patient-Centered Digital Health Records}
Brands et al. (2022) conducted a systematic review of 81 studies with over 1.6 million participants \cite{brands2022}. Findings show PCDHRs improve healthcare utilization (77\% of studies), disease knowledge (70\%), patient engagement and treatment adherence (56\%), and self-management (53\%). Common features include secure messaging, prescription refills, video consultation, and appointment management. Patient satisfaction was positive in 93\% of studies.

\section{Medication Adherence and Digital Health Interventions}
Digital health interventions integrated with electronic health records and patient portals significantly improve medication adherence \cite{brands2022}. Medication tracking, refill reminders, and adherence monitoring are particularly effective for chronic conditions such as diabetes, hypertension, and cardiovascular diseases—conditions increasingly prevalent in Bangladesh.

\section{Progressive Web Applications in Healthcare}
PWAs combine web application reach with native app capabilities. Key advantages include offline functionality, push notifications, and installation without app stores. These features are critical for medication reminders and appointment alerts, especially in areas with unreliable connectivity. PWAs address platform fragmentation and update distribution challenges of traditional mobile health applications.

\section{Integration of AI, EHRs, and Wearables}
Fernandes Prabhu et al. (2025) identified four critical priorities for integrated digital health systems \cite{fernandes2025}: (i) interoperability and data unification from wearables, EHRs, and self-reports; (ii) actionable personalization with predictive insights; (iii) trust and transparency in AI recommendations with clinician oversight; (iv) usability through intuitive interfaces.

\section{Family Health Management and Digital Health Records}
Family health management requires managing information for multiple members, coordinating care across providers, and maintaining comprehensive health histories. Family-centered approaches improve health outcomes, particularly for chronic disease management. Digital health records for family use must address data privacy, access control, and information sharing while maintaining compliance with healthcare regulations applicable in Bangladesh.

\section{Need for Integrated Family Health Management Systems}
MediTrack directly addresses the identified gaps in current digital health solutions by offering:
\begin{itemize}
    \item Unified family health profiles enabling comprehensive tracking across multiple family members.
    \item Integrated medication database with scheduling and adherence monitoring capabilities.
    \item Centralized appointment management reducing missed consultations and improving care coordination.
    \item Secure health report storage and organization facilitating better communication with healthcare providers.
    \item Analytics and insights providing actionable health trend analysis for proactive management.
    \item Offline-capable PWA ensuring accessibility even with unstable internet connectivity.
    \item Modern, responsive design optimized for mobile and desktop devices.
\end{itemize}

% =========================================
% CHAPTER 3: SYSTEM ANALYSIS AND DESIGN
% =========================================
\chapter{System Analysis and Design}

\section{Problem Statement}
Without a proper system, family health management leads to fragmented health records, medication non-adherence, missed appointments, lost medical documents, and difficulty in tracking health trends across family members. The goal of MediTrack is to design and implement a web-based, mobile-friendly, offline-capable system where users can comprehensively manage family health information, track medications, schedule appointments, store health reports, and gain insights through analytics in a unified platform.

\section{Requirements}

\subsection{Functional Requirements}
\textbf{User-side:}
\begin{itemize}
    \item Register, verify email, and log in (with Google OAuth support).
    \item Create and manage family member profiles with health information.
    \item Add, edit, and manage medications in a comprehensive database.
    \item Create medication schedules and track adherence for family members.
    \item Schedule, view, and manage medical appointments.
    \item Upload, organize, and view health reports and medical documents.
    \item View analytics and health insights through interactive charts and visualizations.
    \item Manage purchase lists for medications with cost tracking.
    \item Track health data (vitals, measurements) for family members.
\end{itemize}

\subsection{Non-Functional Requirements}
\begin{itemize}
    \item \textbf{Performance:} Fast loading on low/mid-range devices; optimized bundle size through Vite's build optimization and code splitting.
    \item \textbf{Availability:} Cloud-hosted backend with automatic scaling capabilities.
    \item \textbf{Security:} JWT-based authentication, protected API routes, secure password hashing with bcryptjs, and Google OAuth integration.
    \item \textbf{Reliability:} Real-time data synchronization via React Query plus offline capability through PWA service workers.
    \item \textbf{Usability:} Responsive design supporting mobile and desktop devices with intuitive user interface.
    \item \textbf{Maintainability:} Type-safe database access through Prisma ORM, modular component architecture, and clear separation of concerns.
\end{itemize}

\section{Data Flow Diagram}

\textbf{Level 0 DFD}
\begin{figure}[H]
\centering
\includegraphics[width=0.8\textwidth]{Images/Level 0 DFD Mermaid.png}
\caption{Level 0 DFD of MediTrack}
\label{fig:dfd_level0}
\end{figure}

\textbf{Level 1 DFD}

\begin{figure}[H]
    \centering
    \includegraphics[width=0.8\textwidth]{Images/Level 1 DFD Mermaid.png}
    \caption{Level 1 DFD of MediTrack}
    \label{fig:dfd_level1}
\end{figure}


\section{Database Schema Design (PostgreSQL with Prisma)}
MediTrack uses PostgreSQL as the relational database, managed through Prisma ORM for type-safe database access and migrations.
\begin{itemize}
    \item \textbf{User} - Stores user authentication information (email, password hash, Google ID), display name, and linked family ID for family dashboard access.
    \item \textbf{FamilyMember} - Stores family member profiles with name, relation, age, and last doctor visit information.
    \item \textbf{Medication} - Stores medication database entries with name, strength, price, pack size, and unit information.
    \item \textbf{Plan} - Links family members to medications with dosage, frequency, and pricing information for medication scheduling.
    \item \textbf{Appointment} - Stores appointment details including doctor, specialty, location, date, time, and associated family member.
    \item \textbf{Report} - Manages health reports and documents with title, date, type, category, file URL, and family member association.
    \item \textbf{HealthData} - Tracks health measurements (vitals, lab values) with type, value, unit, and date for family members.
    \item \textbf{PurchaseListItem} - Manages medication purchase lists with quantity, price, and total cost calculations.
    \item \textbf{PurchaseListLog} - Maintains historical records of purchase list completions with summary information.
    \item \textbf{Device} - Tracks user devices for push notification support with device tokens and last active timestamps.
\end{itemize}

\begin{figure}[H]
    \centering
    \includegraphics[width=0.8\textwidth]{Images/Schema Diagram.png}
    \caption{Database Schema Diagram of MediTrack}
    \label{fig:schema}
\end{figure}



% =========================================
% CHAPTER 4: IMPLEMENTATION AND RESULTS
% =========================================
\chapter{Implementation and Results}

\section{Architectural Pattern}
MediTrack uses a layered client–server architecture:
\begin{itemize}
    \item \textbf{React Frontend (Vite):} Modern, fast frontend framework providing component-based UI architecture with optimized builds.
    \item \textbf{Express.js Backend:} RESTful API server handling business logic, authentication, and database operations.
    \item \textbf{Prisma ORM:} Type-safe database access layer providing migrations, query building, and relationship management.
    \item \textbf{Repository Pattern:} Database operations are abstracted through Prisma models, ensuring separation of concerns.
    \item \textbf{Middleware Pattern:} Express middleware handles authentication (JWT), CORS, error handling, and request validation.
    \item \textbf{Observer Pattern:} React Query provides real-time data synchronization and caching, automatically updating UI when data changes.
\end{itemize}

\section{Implementation Results}
The implemented system provides:
\begin{itemize}
    \item \textbf{Family Member Management:} Users can create comprehensive profiles for all family members, track relationships, ages, and health information.
    \item \textbf{Medication Database:} Comprehensive medication tracking with details including name, strength, price, pack size, and unit information.
    \item \textbf{Medication Scheduling:} Plan creation linking family members to medications with dosage, frequency, and cost tracking.
    \item \textbf{Appointment Management:} Centralized scheduling system for tracking medical appointments with doctor, specialty, location, and time information.
    \item \textbf{Health Reports:} Secure document storage and organization with categorization, date tracking, and family member association.
    \item \textbf{Health Data Tracking:} Measurement recording for various health metrics (vitals, lab values) with type, value, unit, and date.
    \item \textbf{Analytics Dashboard:} Interactive visualizations using Recharts providing insights into health trends and medication usage.
    \item \textbf{Purchase List Management:} Medication purchase tracking with quantity, pricing, and historical logging capabilities.
    \item \textbf{Real-time Updates:} React Query ensures data synchronization across components and devices.
    \item \textbf{Offline Capability:} PWA features enable basic functionality even with unstable internet connectivity.
\end{itemize}

\newpage

\LARGE\textbf{Application Screenshots}
\begin{figure}[H]
\centering

% --- Row 1 ---
\begin{subfigure}[b]{0.45\textwidth}
    \centering
    % Image placeholder - user will insert Sign In screenshot
    \caption{Sign In / Sign Up}
    \label{fig:signin}
\end{subfigure}
\hfill
\begin{subfigure}[b]{0.45\textwidth}
    \centering
    % Image placeholder - user will insert Dashboard screenshot
    \caption{Dashboard / Overview}
    \label{fig:dashboard}
\end{subfigure}

\vspace{0.5cm}

% --- Row 2 ---
\begin{subfigure}[b]{0.45\textwidth}
    \centering
    % Image placeholder - user will insert Family Members screenshot
    \caption{Family Members Management}
    \label{fig:family}
\end{subfigure}
\hfill
\begin{subfigure}[b]{0.45\textwidth}
    \centering
    % Image placeholder - user will insert Medication Database screenshot
    \caption{Medication Database}
    \label{fig:medication}
\end{subfigure}
\end{figure}

\newpage

\begin{figure}[H]
\centering

% --- Row 1 ---
\begin{subfigure}[b]{0.45\textwidth}
    \centering
    % Image placeholder - user will insert Appointment Management screenshot
    \caption{Appointment Management}
    \label{fig:appointment}
\end{subfigure}
\hfill
\begin{subfigure}[b]{0.45\textwidth}
    \centering
    % Image placeholder - user will insert Health Reports screenshot
    \caption{Health Reports}
    \label{fig:reports}
\end{subfigure}

\vspace{0.5cm}

% --- Row 2 ---
\begin{subfigure}[b]{0.45\textwidth}
    \centering
    % Image placeholder - user will insert Analytics screenshot
    \caption{Analytics and Insights}
    \label{fig:analytics}
\end{subfigure}
\hfill
\begin{subfigure}[b]{0.45\textwidth}
    \centering
    % Image placeholder - user will insert Purchase List screenshot
    \caption{Purchase List}
    \label{fig:purchaselist}
\end{subfigure}
\end{figure}



% =========================================
% CHAPTER 5: TOOLS AND TECHNOLOGIES
% =========================================
\chapter{Tools and Tech Stack}

\section{Core Framework}
\textbf{React 18} with \textbf{Vite} provides fast development, optimized builds, and component-based architecture.

\section{UI and Styling}
\textbf{Tailwind CSS} provides utility-based styling with minimal CSS bundle. \textbf{Lucide React} offers lightweight, tree-shakeable icons.

\section{Backend and Database}
\textbf{Express.js} provides RESTful API with middleware support. \textbf{Prisma ORM} manages PostgreSQL with type-safe access, migrations, and relationship management.

\section{Authentication and Security}
\textbf{JSON Web Tokens (JWT)} provide stateless authentication. \textbf{Google OAuth} enables Google account authentication. \textbf{bcryptjs} handles secure password hashing.

\section{State Management and Data Fetching}
\textbf{React Query (TanStack Query)} manages server state with automatic caching, background updates, and optimistic updates.

\section{Analytics and Visualization}
\textbf{Recharts} provides interactive, responsive charts for health data visualization.

\section{PWA and Reporting}
Service workers enable offline functionality and PWA installation. \textbf{jsPDF} provides client-side PDF generation for health reports.

\chapter*{Conclusion}
MediTrack provides a unified, offline-capable platform for family health management using React, Vite, Express.js, Prisma, and PWA standards. The system addresses medication adherence, appointment coordination, and health information fragmentation, enabling proactive health management and improved health outcomes.

\bibliographystyle{ieeetr}
\begin{thebibliography}{9}
\bibitem{brands2022}
M. R. Brands, S. C. Gouw, M. Beestrum, R. M. Cronin, K. Fijnvandraat, and S. M. Badawy, ``Patient-Centered Digital Health Records and Their Effects on Health Outcomes: Systematic Review,'' \textit{Journal of Medical Internet Research}, vol. 24, no. 12, p. e43086, Dec. 2022, doi: 10.2196/43086.

\bibitem{fernandes2025}
D. Fernandes Prabhu, V. Gurupur, A. Stone, and E. Trader, ``Integrating Artificial Intelligence, Electronic Health Records, and Wearables for Predictive, Patient-Centered Decision Support in Healthcare,'' \textit{Healthcare}, vol. 13, no. 21, p. 2753, Oct. 2025, doi: 10.3390/healthcare13212753.
\end{thebibliography}

\end{document}